\section{Classification of the second-order rows: Herfurtner's list}\label{sec:herf}

Sections~\ref{sec:sources} and \ref{sec:rootrows} produced two families of
second-order integral rows with apparently unrelated normalisations: Zagier's
$(an^2+an+b,\;cn^2)$ and the root rows'
$(An^2+\tfrac A2n+B,\;C(2n-1)^2)$. This section shows that both shapes are
\emph{forced}, that they are two entries of a list of nine, and that the list is
matched against Herfurtner's classification of rational elliptic surfaces with
four singular fibres \cite{Herfurtner1991}. The outcome is a classification
theorem (Theorem~\ref{thm:herfclass}) whose proved part is the local analysis
and whose verified part is an exhaustive integrality scan.

\subsection{The exponents of a second-order row}

Let $\tht=t\,d/dt$ and let
\begin{equation}\label{eq:R2gen}
(n+1)^2u_{n+1}=P(n)u_n-Q(n)u_{n-1},\qquad u_0=1,\quad u_{-1}=0,
\end{equation}
with $P=an^2+bn+c$ and $Q=dn^2+en+f$ in $\Q[n]$. Then $y=\sum_nu_nt^n$ is
annihilated by
\begin{equation}\label{eq:Lgen}
L=\tht^2-t\,P(\tht)+t^2Q(\tht+1)
 =t^2R(t)\,\partial_t^2+t\,S(t)\,\partial_t+\bigl(-ct+(d{+}e{+}f)t^2\bigr),
\end{equation}
$R(t)=1-at+dt^2$, $S(t)=1-(a+b)t+(3d+e)t^2$.

\begin{lemma}[exponent dictionary]\label{lem:exps}
\Proved{} Assume $d\neq0$ and $a^2\neq4d$, so that $L$ has exactly four singular
points $0,t_1,t_2,\infty$ with $R(t_i)=0$, $t_i=1/\lambda_i$, $\lambda^2-a\lambda+d=0$.
Put $T(t):=S(t)-tR'(t)=1-bt+(d+e)t^2$. Then the local exponents are
\[
(0,0)\ \text{at }0;\qquad
\Bigl(0,\ \rho_i=-\tfrac{T(t_i)}{t_iR'(t_i)}\Bigr)\ \text{at }t_i;\qquad
(1-r_1,\,1-r_2)\ \text{at }\infty,
\]
$r_1,r_2$ being the roots of $Q$; and $\rho_1+\rho_2=r_1+r_2=-e/d$ \textup{(}Fuchs\textup{)}.
The point $t=0$ is a MUM point \textup{(}logarithmic, exponents equal\textup{)}.
\end{lemma}

\begin{proof}
Expanding \eqref{eq:Lgen} gives the displayed $\partial_t$-form. At a simple zero
$t_i$ of $R$ the indicial polynomial is $\rho\bigl[(\rho-1)t_iR'(t_i)+S(t_i)\bigr]$,
whence the exponents $0$ and $1-S(t_i)/(t_iR'(t_i))=-T(t_i)/(t_iR'(t_i))$. At
$\infty$, $y\sim t^{\mu}$ forces $Q(\mu+1)=0$, so the exponents in $s=1/t$ are
$-\mu=1-r_i$. The sum of all six exponents is $2=(4-2)\binom22$ as it must be, and
equals $\rho_1+\rho_2+2-(r_1+r_2)$.
\end{proof}

\begin{theorem}[normalisation classes]\label{thm:norm}
\Proved{} With the hypotheses of Lemma~\ref{lem:exps}, $\rho_1=\rho_2=\rho$ if and
only if
\begin{equation}\label{eq:normclass}
b=(1-\rho)a\qquad\text{and}\qquad e=-2\rho d .
\end{equation}
Writing $Q(n)=C(Mn-j_1)(Mn-j_2)$, so that $r_i=j_i/M$, this reads
\[
\rho=\frac{j_1+j_2}{2M},\qquad
\delta_\infty:=r_2-r_1=\frac{j_2-j_1}{M},\qquad
P(n)=A\Bigl(n^2+\bigl(1-\rho\bigr)n\Bigr)+B ,
\]
with $A=a$, $B=c$ the accessory parameter and $D=CM^2=d$.
\end{theorem}

\begin{proof}
$\rho_1=\rho_2=\rho$ says $T(t_i)+\rho\,t_iR'(t_i)=0$ for $i=1,2$, i.e.\
$R\mid T+\rho tR'$. Both are quadratics with constant term $1$, so they are
equal; comparing coefficients gives $b+\rho a=a$ and $e+2\rho d=0$.
\end{proof}

Two special cases recover the two normalisations of this paper.
\begin{itemize}
\item $\rho=0$ and $\delta_\infty=0$ give $b=a$, $e=f=0$: \emph{Zagier's shape}
$(an^2+an+c,\ dn^2)$. Geometrically $\rho=0$ says both finite singular fibres are
semistable and $\delta_\infty=0$ says the fibre at $\infty$ is too.
\item $\rho=\tfrac12$ and $\delta_\infty=0$ give $b=\tfrac a2$ and
$Q=C(2n-1)^2$: the \emph{root-row shape} of \S\ref{sec:rootrows}, observed
empirically for all nine $\Sym^1$ square roots. Geometrically the two finite
fibres have order-two monodromy ($III$ or $III^*$) and the fibre at $\infty$ has
half-integral exponents ($I_m^*$).
\end{itemize}

\subsection{Kodaira admissibility}

\begin{theorem}\label{thm:kodaira}
\Proved{} Suppose the local system $\ker L$ is, up to a rational gauge factor and
a quadratic character twist, $R^1\pi_*\Q$ of an elliptic surface
$\pi:\mathcal E\to\PP^1_t$. Then
\[
\rho_1,\ \rho_2,\ \delta_\infty\ \in\ \{0,\tfrac12,\tfrac13,\tfrac23\}\pmod 1 .
\]
\end{theorem}

\begin{proof}
Rational gauge and quadratic twist shift exponents by $\tfrac12\Z$ at each point
and do not change exponent \emph{differences} mod $1$. The difference at a point
is $\tfrac1{2\pi i}\log$ of the ratio of the eigenvalues of the local monodromy;
Kodaira's list gives eigenvalue pairs $(1,1)$ for $I_n$, $(-1,-1)$ for $I_n^*$,
$(\pm i)$ for $III,III^*$, $(e^{\pm2\pi i/3})$ for $IV,IV^*$ and
$(e^{\pm\pi i/3})$ for $II,II^*$; the ratios are $1,1,-1,e^{\mp 4\pi i/3},
e^{\mp2\pi i/3}$.
\end{proof}

\begin{corollary}\label{cor:excluded}
\Proved{} The $\Sym^1$ square roots of Cooper's $s_{10}$ and $s_{18}$, whose
trailing coefficients are $-4(8n-3)(8n-5)$ and $12(8n-3)(8n-5)$
\textup{(}so $\delta_\infty=\tfrac14$, an order-eight local monodromy\textup{)},
and the positive-score candidate row of \cite[\S4.4]{ProjNoncong}
\textup{(}$\rho\equiv\tfrac16$\textup{)}, are \emph{not} Picard--Fuchs systems of
an elliptic surface, in any gauge. These are three of the four second-order rows
in the census whose Apéry limits resist identification.
\end{corollary}

\begin{corollary}[Galois rigidity; cf.\ {\cite[Thm.~N3]{ProjNoncong}}]\label{cor:galois}
\Proved{} If $a^2-4d$ is not a square then $t_1,t_2$ are conjugate, hence
$\rho_1=\rho_2$ and the two movable fibres have the same Kodaira class. A
configuration whose movable fibres differ therefore has
$\lambda_1,\lambda_2\in\Z$, so $|\lambda_2|\ge1$ and the row's score
$\log(1/|\lambda_2|)-k$ is $\le-k<0$.
\end{corollary}

\subsection{Herfurtner's list and the cross-ratio invariant}

Herfurtner \cite{Herfurtner1991} classifies the minimal elliptic surfaces over
$\PP^1$ with section, non-constant $\mathcal J$ and exactly four singular
fibres. Since $e(I_0^*)=6$, a rational such surface ($\sum e=12$) has at most one
fibre in $T^-=\{I^*_n,IV^*,III^*,II^*\}$; his Table~3 lists $50$ twist classes and
``transfer of $*$'' yields $56$ fibre quadruples, of which $38$ are rigid (all
fibres in $T^+$; $\mathcal J$ is a Belyi map) and $18$ move in one-parameter
families. The six semistable ones are Beauville's \cite{Beauville1982}.

\begin{definition}\label{def:crossinv}
The row \eqref{eq:R2gen} has four singular points $0,t_1,t_2,\infty$, of which
$0$ is the MUM cusp and $\{t_1,t_2\}$ is an unordered pair. Its
\emph{cross-ratio invariant} is
\[
\mathcal I:=\frac{(1+z)^2}{z}=\frac{(\lambda_1+\lambda_2)^2}{\lambda_1\lambda_2}
=\frac{a^2}{d},\qquad z=\frac{t_1}{t_2} ,
\]
invariant under $z\mapsto1/z$, hence under both allowed relabelings.
\end{definition}

\begin{theorem}[finiteness]\label{thm:finite}
\Proved{} A row \eqref{eq:R2gen} with $a,d\in\Z$ is the Picard--Fuchs system of a
\emph{rigid} Herfurtner configuration only if $\mathcal I=a^2/d$ equals the
cross-ratio invariant of that configuration under an assignment compatible with
Lemma~\ref{lem:exps}. Since $\mathcal I\in\Q$, every configuration whose
invariant is irrational or non-real is excluded outright. Per class the
admissible sets are finite and explicit; for instance
\[
\text{all four fibres semistable:}\quad
\mathcal I\in\Bigl\{0,\ 3,\ \tfrac92,\ -\tfrac{49}8,\ \tfrac{100}9,\
\tfrac{289}{72},\ -121\Bigr\},
\]
\[
\text{root-row class }(\rho,\delta_\infty)=(\tfrac12,0):\quad
\mathcal I\in\Bigl\{\tfrac{484}{125}\ (I_1I_5III\,III),\ 0\ (I_2I_4III\,III),\
-12\ (I_3I_3III\,III)\Bigr\}.
\]
\end{theorem}

Zagier's six realise six of the seven semistable values, in the order
$\mathbf B\,(3)$, $\mathbf E\,(\tfrac92)$, $\mathbf A\,(-\tfrac{49}8)$,
$\mathbf C\,(\tfrac{100}9)$, $\mathbf F\,(\tfrac{289}{72})$,
$\mathbf D\,(-121)$; the discarded eighth value $65+\tfrac{682}{25}\sqrt5$
belongs to $I_5I_5I_1I_1$ with the two $I_1$'s taken as the movable pair, and is
irrational. \emph{The six of Beauville--Zagier is, on the recurrence side, the
number of rational cross-ratios in Herfurtner's four-cusp block.}

\subsection{The $\mathcal J$-map test}

The cross-ratio test ignores the accessory parameter; a sufficient test is
supplied by the nome. Let $q=t\exp(g/y_0)$ be the canonical nome of $L$ at the
MUM point. If $t$ is a Hauptmodul of a genus-zero $\Gamma$ with
$t=c\,q_h+O(q_h^2)$ at a cusp of width $h$ then $q=c\,q_h$, hence
$e^{2\pi i\tau}=(q/c)^h$; so $\Gamma$ is conjugate into $\PSL_2(\Z)$ --- i.e.\
$L$ is projectively the Picard--Fuchs system of an elliptic surface over
$\PP^1_t$ --- iff $J\bigl(\gamma q(t)^h\bigr)$ is a rational function of $t$ for
some $h\in\Z_{>0}$ and some scalar $\gamma=c^{-h}$, and then
$\deg\mathcal J=\deg J(t)$. This is a finite exact linear-algebra computation
(\texttt{lattice/herfurtner/05\_jtest.gp}).

\begin{theorem}[which root rows are elliptic surfaces]\label{thm:whichroots}
\VerifiedR{exact, series order 56, degrees $\le12$, $\gamma$ over $\pm m^{\pm h}$
for $29$ values of $m$} Of the nine $\Sym^1$ square roots of
\cite[\S7--8]{ProjSporadicScan2}, exactly two are Picard--Fuchs systems of a
rational elliptic surface over their own $\PP^1_t$:
\[
\sqrt{\mathrm{AZ}(11,5,125)}\ \leftrightarrow\ I_1I_5\,III\,III\ (\deg\mathcal J=6,\ h=1),
\qquad
\sqrt{\mathrm{AZ}(9,3,-27)}\ \leftrightarrow\ I_3I_3\,III\,III\ (\deg\mathcal J=6,\ h=3),
\]
in agreement with $\mathcal I=\tfrac{484}{125}$ and $\mathcal I=-12$. The
remaining seven --- including Apéry's own square root, i.e.\ Beukers' Theorem~3
row --- have $\mathcal I\notin\{\tfrac{484}{125},0,-12\}$ and fail the
$\mathcal J$-test: their projective monodromy is an Atkin--Lehner quotient
$\Gamma_0(N)+W$, commensurable with $\PSL_2(\Z)$ but not conjugate into it.
\end{theorem}

The mechanism is transparent on Apéry's curve: $t=(\eta_1\eta_6/\eta_2\eta_3)^{12}$
is $W_6$-invariant of degree two on $X_0(6)$, so $\PP^1_t=X_0(6)/W_6$; the
universal elliptic curve does not descend along that quotient, only its
symmetric square does. That is why Apéry's \emph{order-three} row lives happily on
$\PP^1_t$ while its square root does not, and the two fixed points of $W_6$ ---
the Fricke fold point $i/\sqrt6$ of Theorem~B and its companion --- are exactly
the two order-two points, i.e.\ the two exponents $\rho=\tfrac12$.

\subsection{The scan}\label{sec:herfscan}

Theorem~\ref{thm:norm} turns the search for integral rows into $26$ independent
three-parameter searches, one for each Kodaira-admissible pair
$(\rho,\delta_\infty)$ with $r_i=j_i/M\in[-1,2]$. Inside the class $(M;j_1,j_2)$
the free integers are $A$, the accessory parameter $B=u_1$ and the leading
coefficient $C$ of $Q$; \texttt{lattice/herfurtner/02\_hscan.c} sweeps
$1\le A\le3000$, $1\le|B|\le150$, $1\le|C|\le\min(4000,12000/M^2)$ (and a
separate $A=0$ pass), testing integrality division-free on $U_n=u_n(n!)^2$,
\[
U_{n+1}=P(n)U_n-n^2Q(n)U_{n-1},\qquad u_n\in\Z\iff v_p(U_n)\ge2v_p(n!),
\]
prime by prime modulo $p^K$ for every $p\le30$. Survivors are re-verified with
exact rational arithmetic to $n=300$. The sweep is $1.27\times10^6$ hits in
about $75$ minutes on twelve cores.

Two filters precede any claim. The Casoratian
$u_nb_{n+1}-u_{n+1}b_n=\pm\prod_{m\le n}Q(m)/((m+1)!)^2$ vanishes identically as
soon as $Q(n_0)=0$ for an integer $n_0\ge1$ (i.e.\ $M\mid j_i$, $j_i/M\ge1$); the
companion is then a rational multiple of the row, the Ap\'ery limit is rational
and the extension class is split. And $a^2=4d$ gives $\lambda_1=\lambda_2$, so
no second growth rate.

\subsection{Rigid classes and family classes}

The scan of \S\ref{sec:herfscan} splits the $26$ Kodaira-admissible
normalisation classes into two kinds, and the split is Herfurtner's own. A class
whose configurations are all \emph{rigid} produces a handful of integral rows;
a class attached to a \emph{one-parameter family} produces infinitely many. The
cleanest instance is $(\rho;\delta_\infty)=(0;1)$, i.e.\ $Q=C(2n-1)(2n+1)$,
where the fibre at $\infty$ has half-integral exponents (an $I^*_m$): there
\[
P(n)=\tfrac A4(2n+1)^2,\qquad Q(n)=C(2n-1)(2n+1)
\]
is integral for $A\in4\Z$ and $C$ in an arithmetic progression --- one parameter
after the scaling $t\mapsto t/\mu$, matching the one modulus of Herfurtner's
$I_1I_1I_1I_3^*$, $I_1I_1I_2I_2^*$ families. Its $C=4$ sub-family
($\lambda^2-A\lambda+16$, $A=8x$) is the classical Legendre--Pad\'e row with
Ap\'ery limit $-\tfrac1{12}\log\frac{x+1}{x-1}$: $k=1$, score $\to\infty$, and
nothing new.

\subsection{Two new rows}\label{sec:herfnew}

Two integral rows produced by the scan are new to this project and are
Picard--Fuchs systems of Herfurtner surfaces. Both are invisible to Zagier's
normalisation ($\rho\neq0$) \emph{and} to the root-row normalisation
($\delta_\infty\neq0$ resp.\ $\rho\neq\tfrac12$), which is why no earlier scan
found them.

\begin{proposition}\label{prop:newrows}
\VerifiedR{integrality $n\le300$ exact; $\mathcal J$-fit exact}
\begin{enumerate}[label=\textup{(\alph*)},leftmargin=2.4em]
\item $(n+1)^2u_{n+1}=(117n^2+78n+21)u_n-441(3n-1)^2u_{n-1}$, $u_0=1$,
\[u_n=1,\,21,\,693,\,23940,\,734643,\,13697019,\dots\in\Z ,\]
class $(\rho;\delta_\infty)=(\tfrac13;0)$, $\mathcal I=\tfrac{169}{49}$,
$\deg\mathcal J=8$, $h=1$: the surface $I_1I_7\,II\,II$
\textup{(}Herfurtner \#30\textup{)}, projective monodromy $\Gamma_0(7)$.
$\lambda^2-117\lambda+3969$ has discriminant $-2187$, so
$|\lambda_1|=|\lambda_2|=63$ and there is no archimedean limit; $k=2$.
The row is integral only after the scaling $\lambda=3$.
\item $(n+1)^2u_{n+1}=(72n^2+36n+6)u_n-108(4n-1)(4n-3)u_{n-1}$, $u_0=1$,
\[u_n=1,\,6,\,90,\,1140,\,-5850,\,-1265004,\dots\in\Z ,\]
class $(\tfrac12;\tfrac12)$, $\mathcal I=3$, $\deg\mathcal J=3$, $h=3$: the
surface $I_3\,III\,III\,III$ \textup{(}\#45, $\mathcal J=X^3/Y^3$\textup{)},
projective monodromy the index-three congruence group of level three
\textup{(}the kernel of $\PSL_2(\Z)\twoheadrightarrow\Z/3$\textup{)}.
$\lambda^2-72\lambda+1728$ has discriminant $-1728$, $|\lambda_i|=41.57$, no
archimedean limit; $k=2$; integral only after $\lambda=6$.
\end{enumerate}
\end{proposition}

The class $(\tfrac12;\tfrac12)$ deserves a remark: its Fuchsian signature is
$(0;2,2,2;\text{one cusp})$, and \emph{two} groups of that signature occur among
integral rows --- the index-three subgroup of $\PSL_2(\Z)$ just described, and
$\Gamma_0(5)+5$, which has the same covolume but is not conjugate into
$\PSL_2(\Z)$ and carries the level-$5$ Fricke row
$(88n^2+44n+6,\,-4(4n-1)(4n-3))$ of \cite[\S5]{ProjNoncong}. The
$\mathcal J$-map test is exactly what separates them.

\begin{remark}[one surface, two classes]
The three rows $\mathbf A,\mathbf C,\mathbf F$ on $I_6I_3I_2I_1$ are related by the explicit cusp move $t^\sharp=t/(1-\lambda t)$, $F^\sharp=(1-\lambda t)F$, which sends the row $(a,b,d)$ with characteristic roots $\lambda,\mu$ to $(\mu-2\lambda,\,b-\lambda,\,\lambda^2-\lambda\mu)$ \Proved{}. The surface fixes the pure local system; the choice of MUM cusp fixes which Eisenstein class degenerates there: $\mathbf C$ and $\mathbf F$ carry $L(2,\chi_{-3})$ and the same $3$-adic limits up to the factor $\tfrac54$ (Remark~\ref{rem:cuspmove}), while $\mathbf A$ carries $\zeta(2)/4$ and has no $3$-adic limit. The orbits of the six rows under cusp moves are $\{\mathbf A,\mathbf C,\mathbf F\},\{\mathbf B\},\{\mathbf D\},\{\mathbf E\}$ ($\mathbf B,\mathbf D$ have irrational characteristic roots; the third placement of $\mathbf E$ is the degenerate row $u_{2m}=\binom{2m}m^2$).
\end{remark}

\subsection{The classification theorem}

\begin{namedthm}[Classification of second-order rows]\label{thm:herfclass}
Let \eqref{eq:R2gen} have $u_n\in\Z$ for all $n$ \textup{(}after the $\lambda^n$
rescaling of \S\ref{sec:rootrows} if necessary\textup{)}, $\deg P,\deg Q\le2$,
$d\neq0$, $a^2\neq4d$, and $Q(n)\neq0$ for $n\ge1$ \textup{(}non-degenerate
Casoratian\textup{)}.
\begin{enumerate}[label=\textup{(\roman*)},leftmargin=2.4em]
\item \Proved{} The exponents are as in Lemma~\ref{lem:exps}; if $a^2-4d\notin\Q^{\times2}$
then $\rho_1=\rho_2=\rho$ and \eqref{eq:normclass} holds, so $(P,Q)$ lies in one of
the normalisation classes indexed by $(\rho,\delta_\infty)$.
\item \Proved{} If moreover $\ker L$ comes from an elliptic surface over
$\PP^1_t$, then $\rho,\delta_\infty\in\{0,\tfrac12,\tfrac13,\tfrac23\}\bmod1$
\textup{(}Theorem~\ref{thm:kodaira}\textup{)}, the fibre configuration is one of
Herfurtner's $56$, and if it is rigid then $a^2/d$ lies in the explicit finite
set of Theorem~\ref{thm:finite}.
\item \VerifiedR{exhaustive integrality scan of all $26$ Kodaira-admissible
classes, $|a|\le3000$, $|c|\le150$, $|d|\le12000$, integrality tested
prime-by-prime to $n=30$ and re-verified exactly to $n=300$}
in the classes attached to \emph{rigid}
configurations the complete list of primitive, Casoratian-non-degenerate rows is
the $28$ rows of \cite[\S6.1]{ProjHerfurtner}. Exactly ten of them have an
elliptic-surface local system: Zagier's six ($\deg\mathcal J=12$), the two root
rows of Theorem~\ref{thm:whichroots} ($\deg\mathcal J=6$), and the two new rows
of Proposition~\ref{prop:newrows} ($\deg\mathcal J=8$ and $3$). In the classes
attached to \emph{one-parameter} families the integral rows themselves form
infinite families.
\item \VerifiedR{same box; a positive score needs $|\lambda_2|<e^{-k}\le e^{-1}$,
a condition on $(a,d)$ alone, so the pass is complete over all
$1.27\times10^6$ scan hits} Of the $69\,434$ hits in Casoratian-non-degenerate
classes, $1\,374$ have $|\lambda_2|<1$, and \emph{exactly two} have a positive
score with $k\ge2$: Ap\'ery's $\zeta(2)$ row $(11,3,-1)$ at $+0.4061$ and
Beukers' row $(136,10,4)$ at $+0.1392$. Every other positive score has $k=1$ and
belongs to the classical Legendre--Pad\'e log family.
\end{enumerate}
\end{namedthm}

\begin{remark}
What (iii) does \emph{not} settle: the eleven Herfurtner one-parameter families
with a moving $\mathcal J$ are not excluded by Theorem~\ref{thm:finite}
(their cross-ratio is free), only by the scan; and the scan's box, though large,
is a box. The proved part of the classification is entirely local
(Lemma~\ref{lem:exps}, Theorems~\ref{thm:norm}, \ref{thm:kodaira},
\ref{thm:finite}); the global input is Herfurtner's theorem, and the finiteness of
the answer is the finiteness of his list together with the rationality of
$a^2/d$.
\end{remark}
